% !TEX program = xelatex
\documentclass[12pt]{article}
% Compile with XeLaTeX to support Thai Unicode text.
\usepackage{fontspec}
\usepackage{amsmath}
\usepackage{amsfonts}
\usepackage{xcolor}
\usepackage{geometry}
\geometry{a4paper, margin=1in}

% Prefer TH Sarabun New (common Windows Thai font). Scale to match math text.
\IfFontExistsTF{TH Sarabun New}{\setmainfont[Scale=MatchLowercase]{TH Sarabun New}}{%
    \IfFontExistsTF{THSarabunNew}{\setmainfont[Scale=MatchLowercase]{THSarabunNew}}{\setmainfont[Scale=MatchLowercase]{Tahoma}}%
}

% Number only top-level sections (1, 2, 3). No 1.1 for subsections.
\setcounter{secnumdepth}{1}

% --- Bit-table helpers (lecture-style alignment) ---
\setlength{\arraycolsep}{4pt}
\renewcommand{\arraystretch}{1.15}
\newcommand{\bit}[1]{\mathtt{#1}}
\newcommand{\rbit}[1]{\textcolor{red}{\mathtt{#1}}}
\newcommand{\coef}[1]{\mathtt{#1}}
\newcommand{\cblank}{\phantom{\coef{+1}}}
\newcommand{\blank}{\phantom{\bit{0}}}

\begin{document}

\title{Lab 6: Multiplication Algorithms Solution}
\author{CPE223 Computer Architecture}
\date{\today}
\maketitle

\section{Booth's Algorithm (Regular)}
Rules: 
\begin{itemize}
    \item 0 0 or 1 1: Shift Right Arithmetic (ASR)
    \item 0 1: $A = A + M$, then ASR
    \item 1 0: $A = A - M$, then ASR
\end{itemize}

\subsection{พิจารณา $-4 \times 5$}
$M = -4$ (1100), $Q = 5$ (0101)

\[
\begin{array}{cc@{\qquad\Longrightarrow\qquad}c}
\begin{array}{c@{\;}cccc@{\;}l}
\phantom{\times}&\bit{1}&\bit{1}&\bit{0}&\bit{0}&(-4)\\
\times&\bit{0}&\bit{1}&\bit{0}&\bit{1}&(+5)\\
\hline
\end{array}
& &
\begin{array}{*{8}{c}l}
\blank&\blank&\blank&\blank&\bit{1}&\bit{1}&\bit{0}&\bit{0}&\\
\blank&\blank&\blank&\blank&\coef{+1}&\coef{-1}&\coef{+1}&\coef{-1}&\\
\hline
% Shift is shown like the notes: keep the 8-bit row, but leave k trailing blanks for shift k
% shift 0: (+4) = 0000 0100
\rbit{0}&\rbit{0}&\rbit{0}&\rbit{0}&\bit{0}&\bit{1}&\bit{0}&\bit{0}&\\
% shift 1: (-4)<<1 = 1111 1000  (leave 1 blank)
\rbit{1}&\rbit{1}&\rbit{1}&\bit{1}&\bit{1}&\bit{0}&\bit{0}&\blank&\\
% shift 2: (+4)<<2 = 0001 0000  (leave 2 blanks)
\rbit{0}&\rbit{0}&\bit{0}&\bit{1}&\bit{0}&\bit{0}&\blank&\blank&\\
% shift 3: (-4)<<3 = 1110 0000  (leave 3 blanks)
\rbit{1}&\bit{1}&\bit{1}&\bit{0}&\bit{0}&\blank&\blank&\blank&\\
\hline
\bit{1}&\bit{1}&\bit{1}&\bit{0}&\bit{1}&\bit{1}&\bit{0}&\bit{0}&(-20)
\end{array}
\end{array}
\]

\subsection{พิจารณา $5 \times -4$}
$M = 5$ (0101), $Q = -4$ (1100)

\[
\begin{array}{cc@{\qquad\Longrightarrow\qquad}c}
\begin{array}{c@{\;}cccc@{\;}l}
\phantom{\times}&\bit{0}&\bit{1}&\bit{0}&\bit{1}&(+5)\\
\times&\bit{1}&\bit{1}&\bit{0}&\bit{0}&(-4)\\
\hline
\end{array}
& &
\begin{array}{*{8}{c}l}
\blank&\blank&\blank&\blank&\bit{0}&\bit{1}&\bit{0}&\bit{1}&\\
\blank&\blank&\blank&\blank&\coef{0}&\coef{-1}&\coef{0}&\coef{0}&\\
\hline
% 4 rows like the notes (including 0-rows); shifting = trailing blanks
% shift 0: 0000 0000
\rbit{0}&\rbit{0}&\rbit{0}&\rbit{0}&\bit{0}&\bit{0}&\bit{0}&\bit{0}&\\
% shift 1: 0000 0000  (leave 1 blank)
\rbit{0}&\rbit{0}&\rbit{0}&\bit{0}&\bit{0}&\bit{0}&\bit{0}&\blank&\\
% shift 2: (-5)<<2 = 1110 1100  (leave 2 blanks)
\rbit{1}&\rbit{1}&\bit{1}&\bit{0}&\bit{1}&\bit{1}&\blank&\blank&\\
% shift 3: 0000 0000  (leave 3 blanks)
\rbit{0}&\bit{0}&\bit{0}&\bit{0}&\bit{0}&\blank&\blank&\blank&\\
\hline
\bit{1}&\bit{1}&\bit{1}&\bit{0}&\bit{1}&\bit{1}&\bit{0}&\bit{0}&(-20)
\end{array}
\end{array}
\]

\newpage
\section{Bit-pair Recoding (Modified Booth's Algorithm)}

\subsection{พิจารณา $-4 \times 5$}
\[
\begin{array}{cc@{\qquad\Longrightarrow\qquad}c}
\begin{array}{c@{\;}cccc@{\;}l}
\phantom{\times}&\bit{1}&\bit{1}&\bit{0}&\bit{0}&(-4)\\
\times&\bit{0}&\bit{1}&\bit{0}&\bit{1}&(+5)\\
\hline
\end{array}
& &
\begin{array}{*{8}{c}l}
\blank&\blank&\blank&\blank&\bit{1}&\bit{1}&\bit{0}&\bit{0}&\\
% show only the two +1 like the notes (others blank)
\blank&\blank&\blank&\blank&\cblank&\coef{+1}&\cblank&\coef{+1}&\\
\hline
% radix-4 (notes-style): trailing blanks indicate the shift
% shift 0: (-4) = 1111 1100
\rbit{1}&\rbit{1}&\rbit{1}&\rbit{1}&\bit{1}&\bit{1}&\bit{0}&\bit{0}&\\
% shift 2: (-4)<<2 = 1111 0000  (leave 2 blanks)
\rbit{1}&\rbit{1}&\bit{1}&\bit{1}&\bit{0}&\bit{0}&\blank&\blank&\\
\hline
\bit{1}&\bit{1}&\bit{1}&\bit{0}&\bit{1}&\bit{1}&\bit{0}&\bit{0}&(-20)
\end{array}
\end{array}
\]

\subsection{พิจารณา $5 \times -4$}
\[
\begin{array}{cc@{\qquad\Longrightarrow\qquad}c}
\begin{array}{c@{\;}cccc@{\;}l}
\phantom{\times}&\bit{0}&\bit{1}&\bit{0}&\bit{1}&(+5)\\
\times&\bit{1}&\bit{1}&\bit{0}&\bit{0}&(-4)\\
\hline
\end{array}
& &
\begin{array}{*{8}{c}l}
\blank&\blank&\blank&\blank&\bit{0}&\bit{1}&\bit{0}&\bit{1}&\\
% show -1 and the trailing 0 like the notes
\blank&\blank&\blank&\blank&\cblank&\coef{-1}&\cblank&\coef{0}&\\
\hline
% shift 0: 0000 0000
\rbit{0}&\rbit{0}&\rbit{0}&\rbit{0}&\bit{0}&\bit{0}&\bit{0}&\bit{0}&\\
% shift 2: (-5)<<2 = 1110 1100  (leave 2 blanks)
\rbit{1}&\rbit{1}&\bit{1}&\bit{0}&\bit{1}&\bit{1}&\blank&\blank&\\
\hline
\bit{1}&\bit{1}&\bit{1}&\bit{0}&\bit{1}&\bit{1}&\bit{0}&\bit{0}&(-20)
\end{array}
\end{array}
\]

\section{Discussion}
\sloppy
\textbf{1) Booth's Algorithm (Regular)}

ประสิทธิภาพของ Booth ขึ้นอยู่กับรูปแบบบิตของ \textit{ตัวคูณ} $Q$ เพราะอัลกอริทึมดูคู่บิต
$(Q_i,Q_{i-1})$ แล้วตัดสินใจว่าในคาบนั้นจะ \textit{บวก/ลบ} $M$ หรือ \textit{เลื่อน (ASR)} อย่างเดียว
ดังนั้นจำนวนครั้งที่ต้องทำ ``บวก/ลบ'' จะขึ้นกับจำนวนครั้งที่บิตของ $Q$ เปลี่ยนสถานะ
(ขอบเขต 0\,$\leftrightarrow$\,1) ในการไล่บิต
\begin{itemize}
    \item กรณี $-4\times 5$ (ตัวคูณ $Q=0101$): บิตสลับ 0/1 ทำให้คู่บิตสลับเป็น
    01 และ 10 เกือบทุกคาบ จึงเกิดการดำเนินการแบบ $+M$ / $-M$ มากที่สุด
    (สำหรับ 4 บิตคือ 4 ครั้ง) ถือเป็น \textit{Worst Case} และแทบไม่ต่างจากการคูณแบบปกติ
    \item กรณี $5\times -4$ (ตัวคูณ $Q=1100$): มีช่วงบิตซ้ำ (11 และ 00)
    หลายคาบจึงเป็นกรณี 11/00 ที่ทำแค่ ASR โดยไม่ต้องบวก/ลบ
    ในตัวอย่างนี้จะมีการทำบวก/ลบเพียง 1 ครั้ง (ที่คู่บิต 10)
    ถือเป็น \textit{Best Case}
\end{itemize}

\textbf{ข้อจำกัดของ Booth:} หากรูปแบบบิตของตัวคูณเข้าใกล้ Worst Case (เช่น 0101)
ประโยชน์ด้านความเร็วจะลดลงมาก

\medskip
\textbf{2) Bit-pair Recoding Multiplication (Radix-4 Modified Booth)}

Bit-pair recoding จะพิจารณา 3 บิต $(Q_{2i+1},Q_{2i},Q_{2i-1})$ แล้วเข้ารหัสเป็นค่าสัมประสิทธิ์ในชุด
$\{-2,-1,0,+1,+2\}$ ทำให้จำนวน partial products ลดลงจากประมาณ $n$ แถว
เหลือประมาณครึ่งหนึ่งคือ $\lceil n/2 \rceil$ แถว (สำหรับตัวคูณ $n$ บิต)
\begin{itemize}
    \item สำหรับแลปนี้ $n=4$ จึงเหลือ 2 แถวในทั้งสองกรณี ($-4\times 5$ และ $5\times -4$)
    \item เมื่อจำนวนแถวลดลง จำนวนรอบ/จำนวนการบวกสะสมก็ลดลงด้วย (เช่น 4 รอบเหลือ 2 รอบ)
    และให้เวลาคำนวณที่สม่ำเสมอกว่า Booth แบบปกติ
\end{itemize}

\fussy

\end{document}